\begin{table*}[htbp]
\centering
\small
\caption{Три предрегистрированные mixed-effects модели: идентифицируемость и статус}
\label{tab:8}
\begin{tabular}{>{\RaggedRight\hspace{0pt}\arraybackslash}p{0.120\linewidth}>{\RaggedRight\hspace{0pt}\arraybackslash}p{0.120\linewidth}>{\RaggedRight\hspace{0pt}\arraybackslash}p{0.120\linewidth}>{\RaggedRight\hspace{0pt}\arraybackslash}p{0.120\linewidth}>{\RaggedRight\hspace{0pt}\arraybackslash}p{0.120\linewidth}>{\RaggedRight\hspace{0pt}\arraybackslash}p{0.120\linewidth}>{\RaggedRight\hspace{0pt}\arraybackslash}p{0.120\linewidth}}
\toprule
Модель & Корпус & Первичный исход & Наблюдений & Ранг матрицы плана & Слитых уровней & Статус \\
\midrule
M1 & машинный & O1 & 1079 & 22 из 23 & 1 & не идентифицируема; подгонка запущена и не сошлась \\
M2 & человеческий & O2 & 803 & 7 из 22 & 15 & не идентифицируема \\
M3 & объединённый & O3 & 1882 & 11 из 11 & 0 & идентифицируема \\
\bottomrule
\end{tabular}
\begin{minipage}{\linewidth}\footnotesize\vspace{2pt}
\textit{Примечание.} Единица анализа — наблюдение — документ. Знаменатель: 1079 машинных, 803 человеческих и 1882 объединённых документа. Данные: серия v2, диагностика идентифицируемости от 2026-07-29. Ранг матрицы плана против числа кандидатных колонок; слитые уровни — факторы, не отличимые по дизайну. Оценок эффектов таблица не приводит: ни одна модель не дала первичный исход. Шкала: шкалы нет, таблица описывает статус процедуры. Сокращения: REML — restricted maximum likelihood; слитый уровень — уровень фактора, линейно зависимый от остальных колонок плана. Пропуски: колонка оценок отсутствует намеренно: M1 не сошлась, M2 неидентифицируема по дизайну, у M3 нет поддержки зарегистрированного estimand, поэтому чисел, которые можно было бы опубликовать, нет. Поправка на множественные сравнения не применялась.
\end{minipage}
\end{table*}
